\scanpage{101}
\mspage{50}

\deleted{\unsure{ce}}\inserted{l'interprétation} qui nous donne aussitôt l'application canonique de degré 2

\[
\begin{aligned}
 \unsure{\vec\Pi|\overline\Pi}&\xrightarrow[\text{degré 2}]{}\Pi\\
 \{i,\underset{\substack{\rotatebox{90}{$\in$}\\\deleted{\unsure{\operatorname{Or}}}\unsure{\mu|\chi}(i)}}e\}&\longmapsto i
\end{aligned}\tag*{(69)}
\]

Il faut \deleted{définir}\inserted{\unsure{définir|décrire} combinatoirement} une opération de $\Lambda^*$ sur $\vec\Pi$, de telle façon que (69) soit compatible avec l'application canonique

\[
\deleted{\unsure{(70)}}\qquad\Lambda^*\longrightarrow Z'_\Lambda=\Lambda^*/\{\pm1_\Lambda\}
\]

Soit donc $(i,e\in\deleted{\operatorname{Or}}\mu(i)\deleted{\unsure{P_{M,i}}})\in\vec\Pi$, $\lambda\in\Lambda^*$, $\unsure{\acute\lambda|\dot\lambda}$ son image ds $Z'_\Lambda$ on veut définir

\[
\left\{\begin{aligned}
 &\lambda(i,e)=\deleted{\unsure{\lambda}}(i',e')\qquad e'\in\operatorname{Or}\deleted{\unsure{P_{M,i'}|P_{i'}}}\mu(i')\\
 &i'=\dot\lambda.i,\quad e'=\mathord{?}
\end{aligned}\right.\tag*{(70)}
\]

\deleted{Or considérons l'application (61) ($s\in S$ choisi)}

\[
\deleted{\operatorname{Et}_{P_{i_0}}(s)\longrightarrow\operatorname{Et}_{P_i}(s)}
\]

\deleted{le choix de $e$ définit une orientation du \unsure{$n$}-polygône $\operatorname{Et}_{P_{i_0}}(s)$, qui devient ainsi un torseur sous $\mathbb Z/n\mathbb Z$ ; pour \deleted{\unsure{tt}}\inserted{le choix d'une} orientation $e'$ de $P_{i'}$, faisant de $\operatorname{Et}_{P_i}$}

Choisissons $s\in S$, alors

\[
\lambda_{\Gamma,s}\in\Lambda\wedge_{\mu_2}(\mu(i)\wedge\mu(i'))\xrightarrow{\text{degré 2}} Z'_\Lambda\text{ --}
\]

est au dessus de $\dot\lambda$ ; pour tout $e'\in\mu(i')$,
% archive sheet number: 100
